
\[
\begin{tikzcd}[column sep=large,row sep=large]
K^\bullet(X_{\bar t}) \arrow[r,"c^i"] \arrow[d,"\sigma"'] &
H^{2i}(X_{\bar t},\mu_n^{\otimes i})\\
K^\bullet(X_{\bar s}) \arrow[r,"c^i"] &
H^{2i}(X_{\bar s},\mu_n^{\otimes i}) \arrow[u,"\sigma^{-1}"]
\end{tikzcd}
\;,
\]
où \(X_{\bar t}\), \(X_{\bar s}\)
sont les fibres \emph{géométriques} de \(X\). D'ailleurs, on sait que maintenant l'homomorphisme \(\sigma^{-1}\) de spécialisation pour la cohomologie est un isomorphisme (SGA 4 XVI 2.2), donc la commutativité du carré précédent peut s'exprimer encore par celle du carré qui s'en déduit en inversant la flèche précédente. Passant alors à la limite projective sur \(n\) dans les carrés commutatifs ainsi obtenus, on trouve un carré commutatif
\begin{equation}\tag{7.13.4}
\begin{tikzcd}[column sep=large,row sep=large]
K^\bullet(X_{\bar t}) \arrow[r] \arrow[d,"\sigma"'] &
H^{2i}(X_{\bar t},\mathbb Z_\ell(i)) \arrow[d,"\wr"']\\
K^\bullet(X_{\bar s}) \arrow[r] &
H^{2i}(X_{\bar s},\mathbb Z_\ell(i)).
\end{tikzcd}
\end{equation}

\textbf{7.13.5.} Une autre façon d'exprimer cette commutativité est de dire que, si \(x_t\) est un élément de \(K^\bullet(X_t)\) et \(x_s\) son image dans \(K^\bullet(X_s)\), alors \(c^i(x_{\bar t})\) et \(c^i(x_{\bar s})\) sont les valeurs, en les points géométriques \(\bar t\) et \(\bar s\) respectivement, d'une \emph{section} de \(R^{2i}f_*(\mathbb Z_\ell(i))\). Cette section s'interprète d'ailleurs aisément, en prenant un élément \(x\in K^\bullet(X)\) qui induit \(x_t\), d'où un \(c^i(x)\in H^{2i}(X,\mathbb Z_\ell(i))\), qui définit la section cherchée de \(R^{2i}f_*(\mathbb Z_\ell(i))\).

\textbf{7.13.6.} Soit \(T\) un schéma. Alors par passage aux gradués associés, l'homomorphisme (7.13.0) définit des homomorphismes
\[
  \psi_i:\Gr^i(T)\longrightarrow H^{2i}(T,\mathbb Z_\ell(i)),
\]
l'image par \(\psi_i\) d'un élément \(x\) de \(\Filt^iK^\bullet(T)\) étant par définition \(c^i(x)\). (Le fait que (7.13.0) est compatible avec la \(\lambda\)-filtration sur 
